% Generated draft: requires mathematical and source review.
\documentclass[11pt]{article}
\usepackage[margin=25mm]{geometry}
\usepackage{fontspec}
\setmainfont{Latin Modern Roman}
\usepackage{amsmath,amssymb,mathtools,unicode-math}
\setmathfont{Latin Modern Math}
\usepackage{enumitem,needspace,xurl,hyperref}
\setlength{\parindent}{0pt}
\setlength{\parskip}{5pt}
\setlength{\emergencystretch}{3em}
\newcommand{\N}{\mathbb{N}}
\newcommand{\Z}{\mathbb{Z}}
\newcommand{\Q}{\mathbb{Q}}
\newcommand{\R}{\mathbb{R}}
\newcommand{\C}{\mathbb{C}}
\newcommand{\F}{\mathbb{F}}
\newcommand{\A}{\mathbb{A}}
\newcommand{\PP}{\mathbb{P}}
\newcommand{\E}{\mathbb{E}}
\newcommand{\eps}{\varepsilon}
\newcommand{\dd}{\,\mathrm{d}}
\begin{document}
\section*{UnsolvedMath ID: 30000866}
\textbf{Source problem code:} catalogue entry
\subsection{Definitions and mathematical statement}
Let \(N\geq1\) be odd and squarefree, let \(K\) be a totally real number field of degree \(d=[K:\Q]\), and let \(\mathcal O=\mathcal O_K\). Write \(K(N)\) for the genus-two paramodular group of level \(N\), as in the source; this group notation is distinct from the coefficient field \(K\). The proposed paramodular correspondence relates weight-two, degree-two paramodular Hecke newforms \(f\), whose Hecke eigenvalues generate \(K\), to isogeny classes of simple abelian varieties \(A/\Q\) satisfying
\[
\dim A=2d,\qquad \operatorname{cond}(A)=N^d,
\qquad A\text{ has endomorphisms by }\mathcal O.
\]
The correspondence is required to satisfy
\[
L(A,s)=\prod_{\sigma:K\hookrightarrow\C}L_{\mathrm{spin}}(f^\sigma,s),
\]
where \(f^\sigma\) denotes the Hecke-eigenvalue conjugate and \(L(A,s)\) is the Hasse--Weil \(L\)-function. The source asks for a bijection, not just existence of an abelian variety for each form.

The raw record additionally states the prime-level specialization: rational nonlift Hecke eigenforms in \(S_2^2(K(p))\) should correspond bijectively to rational abelian surfaces of conductor \(p\), with
\[
L_{\mathrm{spin}}(f,s)=L(A,s).
\]
Here rational means rational Hecke eigenvalues on the form side and definition over \(\Q\) on the abelian-variety side. Odd prime \(p\) is compatible with the general hypotheses; the raw specialization itself says only ``prime.''

Several choices affecting a precise bijection are not supplied. The general formulation does not explicitly impose the nonlift restriction that appears in the prime specialization. It also does not say whether forms are normalized eigenforms, eigenlines, or Galois orbits; the product over conjugates naturally loses the distinction between conjugate forms, so no orbit convention is silently chosen. Likewise, ``simple,'' the field of definition of the \(\mathcal O\)-endomorphisms, whether this action is part of the isogeny data, and the exact local \(L\)-factor normalization require source review.
\subsection{Short English statement}
The proposed paramodular bijection pairs suitable weight-two genus-two newforms with simple abelian varieties over \(\Q\) of dimension \(2d\) and conductor \(N^d\), matching the product of conjugate spinor \(L\)-functions with the Hasse--Weil function. The raw source also requires nonlift forms in its rational prime-level specialization.
\subsection{Sources}
\begin{itemize}
\item Local UnsolvedMath registry, permanent ID \(30000866\), source code catalogue entry, ``Paramodular Conjecture for Abelian Varieties.'' No individual proposer is supplied.
\item \emph{Modulformen} (2008), Oberwolfach Reports, DOI 10.4171/owr/2007/51, pp. 3013--3072: \url{https://doi.org/10.4171/owr/2007/51}.
\item Supplied reference, ``Brumer--Kramer, Paramodular Abelian Varieties of Odd Conductor'': \url{https://arxiv.org/abs/1004.4699}.
\item Supplied reference, ``paramodular conjecture overview'': \url{https://www.lmfdb.org/knowledge/show/g2c.paramodular_conjecture}.
\item Supplied literature triage, checked 2026-08-21: status ``partially solved.'' It reports substantial computational and conditional progress for rational abelian surfaces, including prime-conductor results, while saying the asserted general bijection is not proved. This dated assessment derives from supplied OWR curation; time-sensitive status and cited claims require independent verification.
\end{itemize}
\end{document}
